\chapter{Considerações Finais}
\label{cap:Consideracoes}

%Falar sobre injeção de prompt!!!!!
%Falar de perguntas inadequadas e como tratar (testar na fase de teste com humanos)

Este trabalho apresentou o desenvolvimento de um sistema de busca semântica baseado na abordagem \textit{Retrieval-Augmented Generation} (\ac{RAG}), voltado para a recuperação e síntese de informações textuais presentes em documentos da Assembleia Legislativa do Estado do Rio Grande do Norte (\ac{Alern}). A solução proposta combina técnicas de Processamento de Linguagem Natural (\ac{PLN}), bancos vetoriais e \textit{Large Language Models} (\ac{LLM}s), culminando na construção de um assistente de busca capaz de interpretar a intenção do usuário e gerar respostas contextualizadas. O sistema busca responder à necessidade de recuperação eficiente de informações em um conjunto textual de grande volume e complexidade, como os encontrados na \ac{Alern}, promovendo maior acessibilidade e transparência no processo legislativo.

Os resultados obtidos demonstraram a viabilidade da abordagem proposta. A combinação entre a recuperação eficiente de fragmentos de texto relevantes, realizada por meio de busca semântica em bancos vetoriais, e a geração de respostas sintéticas pelo \ac{LLM}, mostrou-se promissora tanto em termos de precisão quanto de completude das informações fornecidas. As avaliações automatizadas realizadas no decorrer dos experimentos confirmaram a capacidade do sistema em identificar documentos relevantes e produzir respostas coesas e alinhadas com as consultas realizadas. A integração com o Ollama, utilizando uma instância do modelo Llama 3.1, em um ambiente de execução com GPU T4 de 16 GB no Google Colab, permitiu superar limitações de infraestrutura, garantindo a execução eficiente dos testes propostos.

Apesar dos avanços obtidos, é importante destacar as limitações encontradas durante o desenvolvimento e avaliação da solução. A precisão dos resultados depende diretamente da qualidade das representações vetoriais geradas e do corpus de documentos disponibilizado. Além disso, a abordagem \ac{RAG}, embora eficaz na mitigação de alucinações, ainda pode apresentar respostas imprecisas em contextos ambíguos ou com dados insuficientes. Diante dessas limitações, sugere-se como trabalhos futuros o aprimoramento do modelo de recuperação, com a incorporação de técnicas de re-ranking, e a exploração de \ac{LLM}s mais robustos, além da ampliação do corpus documental com dados atualizados e enriquecidos semanticamente.

Por fim, este estudo contribui para o avanço das pesquisas na área de busca semântica e aplicação de \ac{LLM}s em contextos institucionais, oferecendo uma solução prática e escalável que atende às demandas de transparência e eficiência no acesso a informações públicas. A implementação do sistema na \ac{Alern} tem o potencial de fortalecer a interação entre a população e o poder legislativo, promovendo o exercício da cidadania por meio do acesso simplificado e inteligente às informações de interesse público.